\scanpage{27}
\mspage{\handwritten{19}}
% Source viewed upright after a 180-degree rotation.
% The enclosing TikZ reconstruction carries the page-wide cancellation strokes.
% Ordinary local cancellations remain \deleted; the whole draft is cancelled
% by the drawn oblique strokes, rather than by an invented horizontal strike.
% The left marginal oval is placed in the same coordinate system as its line.
\reversemarginpar
\figure{%
\begingroup
\csname tagsleft@true\endcsname
\begin{tikzpicture}[x=.01cm,y=-.01cm,scale=.74,transform shape,
  every node/.style={anchor=north west,inner sep=0pt,align=left,
    font=\fontsize{10}{20}\selectfont\ttfamily}]
\path[use as bounding box] (55,180) rectangle (1650,2145);

\node at (230,202) {\begin{minipage}{14.1cm}
2.1.1.\quad On peut dire encore qu'une structure de biextension sur le $G_{P\times Q}$-torseur $E$ est définie par deux lois de composition non partout définies $\handwritten{\phi}(x,y)$ et $\handwritten{\psi}(x,y)$ sur $E$, la première étant définie lorsque \deleted{$j($} $\mathrm{pr}_2j(x)=\mathrm{pr}_2j(y)$, la deuxième lorsque $\mathrm{pr}_1j(x)=\mathrm{pr}_1j(y)$ \deleted{(où} (où $j:E\handwritten{\longrightarrow}P\times Q$ est le morphisme structural du torseur). Les axiomes, en plus des associativités du type
\end{minipage}};

\node[anchor=center] at (107,374) {$\handwritten{\unsure{\psi|\Psi}}$};
\draw[line width=.45pt] (116,324) .. controls (87,324) and (58,389) .. (74,424)
  .. controls (91,454) and (145,391) .. (139,351)
  .. controls (138,339) and (128,329) .. (116,324);

% Two small isolated ink dots at the left of the abandoned formula.
\fill (237,633) circle[radius=.014cm];
\fill (233,640) circle[radius=.014cm];

\node at (465,616) {\begin{minipage}{8.35cm}
\[
\deleted{\handwritten{\phi}(gx,y)\unsure{=}g\handwritten{\phi}(x,y),\quad\handwritten{\phi}(x,gy)\unsure{=}g\handwritten{\phi}(x,y),\quad\handwritten{\phi}(}
\]
\end{minipage}};

\node at (230,666) {\begin{minipage}{14.1cm}
\[
\handwritten{\phi}(gx,y)=\handwritten{\phi}(x,gy)=g\handwritten{\phi}(x,y)\qquad(x,y\handwritten{\in}E(S),\ g\handwritten{\in}G(S))
\tag*{(2.1.1.1)}
\]
\end{minipage}};

\node at (230,751) {\begin{minipage}{14.1cm}
et des relations analogues en $\handwritten{\psi}$, s'expriment alors par l\deleted{a}'associativité et la commutativité de \unsure{\deleted{s}c|\deleted{c}s}es lois de composition partielles, et par la relation
\end{minipage}};

\node at (230,946) {\begin{minipage}{14.1cm}
\[
\handwritten{\phi}(\handwritten{\psi}(x,y),\handwritten{\psi}(z,t))=\handwritten{\psi}(\handwritten{\phi}(x,z),\handwritten{\phi}(y,t))
\tag*{\raisebox{10pt}{(2.1.1.2)}}
\]
\end{minipage}};
% The two local hooked strokes are distinct from the page-wide cancellation.
\draw[line width=.45pt] (301,1005) .. controls (281,982) and (436,953) .. (562,941) -- (565,928);
\draw[line width=.6pt] (441,934) .. controls (465,950) and (421,966) .. (399,982)
  .. controls (382,994) and (399,1010) .. (424,1018)
  .. controls (472,1022) and (522,1018) .. (571,1021) -- (574,1034);

\node at (230,1039) {\begin{minipage}{14.1cm}
pour $x,y,z,t$ respectivement dans $E_{p,q}(S),E_{p,q'}(S),E_{p',q}(S),E_{p',q'}(S)$.
\end{minipage}};

\node at (350,1088) {\inserted{\handwritten{\uline{Sections unité}.}}};
\draw[line width=.45pt] (652,1107) .. controls (659,1134) and (601,1135) .. (444,1136)
  .. controls (423,1134) and (412,1151) .. (418,1167);

\node at (230,1126) {\begin{minipage}{14.1cm}
2.2.\quad Il y a lieu de considérer la section unité de $\mathbb{E}\deleted{t}$ considéré comme extension de $Q_P$ par $G_P$,
\end{minipage}};

\node at (230,1252) {\begin{minipage}{14.1cm}
\[
e_{E/P}\handwritten{\in\Gamma}(E/P)=\mathrm{Hom}_P(P,E)\ ,
\tag*{(2.2.1)}
\]
\end{minipage}};

\node at (230,1334) {\begin{minipage}{14.1cm}
et de façon symétrique la section unité de $E$\handwritten{,} considéré comme extension de $P_Q$ par $G_Q$,
\end{minipage}};

\node at (230,1462) {\begin{minipage}{14.1cm}
\[
e_{E/Q}\handwritten{\in\Gamma}(E/Q)=\mathrm{Hom}_Q(Q,E)\ .
\tag*{(2.2.2)}
\]
\end{minipage}};

% This is one paragraph. Its two interlinear additions are placed separately
% at their source positions so that their relation to the cancellation is kept.
\node at (230,1547) {\begin{minipage}{14.1cm}
En particulier, \uline{le torseur} $E$ \uline{au dessus de} $P\times Q$ (de \unsure{\deleted{g}G|\deleted{G}g}roupe $G_{P\times Q}$) \uline{est canoniquement trivialisé au dessus de} $1_P\times Q$ \uline{par} $e_{E/P}$, \uline{et au dessus de} $P\times1_Q$ \uline{par} $e_{E/Q}$. \uline{\unsure{\deleted{C}c|\deleted{c}C}es deux trivialisations coïncident au dessus de la section unité} $1_P\times1_Q$ \uline{de} $P\times Q$, et définissent une même section
\end{minipage}};
\node at (590,1654) {\inserted{De plus,}};
\node at (591,1676) {\deleted{\handwritten{\illegible{1}}}};
\draw[line width=.6pt] (767,1668) .. controls (780,1699) and (645,1682) .. (600,1687)
  .. controls (569,1690) and (571,1715) .. (578,1734);
\node at (1071,1728) {\inserted{\handwritten{donc}}};
\draw[line width=.45pt] (1153,1738) .. controls (1169,1756) and (1082,1752) .. (1085,1781);

\node at (230,1819) {\begin{minipage}{14.1cm}
\[
e_{E\deleted{/S}}\handwritten{\in\Gamma}(E\handwritten{)}\deleted{\unsure{/S)|/P)|/Q)}}\ ,
\tag*{(2.2.3)}
\]
\end{minipage}};

\node at (230,1898) {\begin{minipage}{14.1cm}
qu'on appellera par abus de langage la \uline{section unité} (ou \uline{section neutre}) de $E$. Pour prouver l'égalité annoncée\deleted{, et une conséquence} des deux sections $e,e'$ de $E$ déduites respectivement de $e_{E/P}$ et $e_{E/Q}$, on note
\end{minipage}};

% Page-wide source strokes; upright coordinates in the original 1685 x 2159 scan.
\draw[line width=.65pt] (1495,351) .. controls (1230,570) and (455,1360) .. (132,1740);
\draw[line width=.55pt] (1555,998) .. controls (1270,1275) and (765,1828) .. (522,2120);
\draw[line width=.5pt] (1558,1427) .. controls (1613,1510) and (1290,1933) .. (1118,2110);
\draw[line width=.45pt] (495,188) .. controls (585,317) and (1005,1260) .. (1235,1892);
\draw[line width=.4pt] (938,238) .. controls (1056,501) and (1400,1065) .. (1485,1420);
\draw[line width=.4pt] (220,623) .. controls (218,864) and (445,1770) .. (435,2140);
\end{tikzpicture}
\endgroup
}
\normalmarginpar
% archive sheet number: 26 (oppositely oriented to the typescript)
